\documentclass[11pt]{article}
\usepackage[a4paper,margin=27mm]{geometry}
\usepackage[T1]{fontenc}
\usepackage{lmodern,amsmath,amssymb,amsthm,mathtools,booktabs,xurl}
\usepackage[hidelinks]{hyperref}
\emergencystretch=3em
\newtheorem{theorem}{Theorem}[section]
\newtheorem{lemma}[theorem]{Lemma}
\newtheorem{corollary}[theorem]{Corollary}
\newtheorem{proposition}[theorem]{Proposition}
\newcommand{\conv}{\operatorname{conv}}
\newcommand{\vol}{\operatorname{vol}}
\newcommand{\BM}{d_{\mathrm{BM}}}
\newcommand{\ee}{\mathfrak e}
\newcommand{\E}{\mathbb E}
\newcommand{\R}{\mathbb R}
\title{Sharp-exponent simplex stability with a quadratic dimension constant}
\author{mxym/math research project\\AI-assisted research manuscript}
\date{7 October 2026}
\begin{document}
\maketitle
\begin{abstract}
For the lower-end projection-pyramid deficit, we prove centroid containment by
every prescribed maximum-volume inscribed simplex with sharp exponent
$1/(d-1)$ and an explicit coefficient at most $4096d^2$ in every dimension
$d\ge3$. The proof combines a conditioning-free, volume-weighted anchor
theorem with an intrinsic first-moment conversion and a relative projection-cap
argument. A truncated simplex gives the sharp obstruction, including the
classification of all its maximum-volume inscribed simplices. The optimal
coefficient lies between linear and quadratic order in the dimension; this
paper does not close that gap. The traditional proof is complete. Existing
Lean releases cover the original sharp-exponent upper theorem and the exact
truncation obstruction; the improved quadratic coefficient is a written
analytic refinement, not represented as fully formalized.
\end{abstract}

\section{Definitions and the main statement}
Let $K\subset\R^d$ be a full-dimensional compact convex body. Its projection
body has support function $h_{\Pi K}(u)=\|u\|_2\vol_{d-1}(K|u^\perp)$.
Write $R(K)=\vol_d(\Pi K)/\vol_d(K)^{d-1}$, and let $\mathcal P K$ be a
pyramid in dimension $d+1$ with base $K$. Set
\[
a(K)=\left(\frac d{d+1}\right)^d\frac{R(\mathcal P K)}{R(K)}-1,
\qquad \ee(K)=a(K)-\frac1{d+1}.
\]
The height and apex location do not affect this affine invariant. For an
inscribed simplex $S$ of maximum volume, let $g_S$ denote its original centroid
and define
\[
E(K,S)=\inf\{t\ge0:K\subset g_S+(1+t)(S-g_S)\}.
\]
Thus $E$ is excess dilation, not the dilation factor itself. For a reference
simplex $\Delta_d$ the affine Banach--Mazur distance to its affine class is
\[
\BM(K,\Delta_d)=\inf\{\lambda\ge1:S'\subset K\subset v+\lambda S',
\ S'\text{ a full-dimensional simplex},\ v\in\R^d\}. \tag{1}
\]
The free translation and optimization in (1) distinguish it from $E(K,S)$.

\begin{theorem}\label{thm:main}
For every $d\ge3$, every such $K$, and every prescribed maximum-volume
inscribed simplex $S$,
\[
\begin{gathered}
\ee(K)\ge0,\qquad E(K,S)\le G_d\ee(K)^{1/(d-1)},\\
G_d=16(d+1)^2[3d^2(d+1)^3(d+2)]^{1/(d-1)}\le4096d^2.
\end{gathered}
\]
The exponent cannot be increased, even for the best choice of maximum $S$
or for $\BM(K,\Delta_d)-1$. Any coefficient valid for every maximum $S$
at this exponent is at least
$(d+1)[(d+1)^2/(d(d-1))]^{1/(d-1)}$.
\end{theorem}
No product, join, smoothness or strict-convexity assumption is imposed on $K$.
All maximum simplices exist by compactness. Replacement of one vertex of a
maximum $S$ by any $x\in K$ shows that every $S$-barycentric coordinate
$\alpha_i(x)$ has $|\alpha_i(x)|\le1$. Consequently $E(K,S)\le d+1$.

\section{Cone law and the exact deficit identity}
Translate an interior point of $K$ to zero, write $v=\vol_d K$, and denote its
surface-area measure by $S_K$. Push the probability measure
$h_K(u)dS_K(u)/(dv)$ on the unit sphere forward by $u\mapsto u/h_K(u)$;
call the resulting law $\nu$. It is centered, bounded and supported on
$\partial K^\circ$, with full affine span. Centering follows from
$\int u\,dS_K(u)=0$; the total mass follows from the divergence formula.
The support of $S_K$ spans $\R^d$ for a full-dimensional body. Since $\nu$ is
centered, linear and affine full span are equivalent. For independent samples
put
\[
A=\E|\det(X_1,\ldots,X_d)|,\quad
B=\E|\det((X_0,1),\ldots,(X_d,1))|,\quad D=B-A.
\]
Then $B>0$ and
\[
a(K)=\frac B{(d+1)A},\qquad
\frac DB=\frac{(d+1)\ee(K)}{1+(d+1)\ee(K)}.
\]
Here is a direct derivation, including the normalization. For a polytope let
$u_i$ be facet area times outward unit normal and $b_i$ facet area times support
number. Set $H=\sum_{|I|=d}|\det(u_i:i\in I)|$ and
$L=\sum_{|I|=d+1}|\det((u_i,b_i):i\in I)|$. Cauchy's formula and the
zonotope volume formula give $\vol\Pi K=H$. For the pyramid with height $t$
and apex above zero, the facet area vectors are
$(0,-v)$ and $(t u_i/d,b_i/d)$, and its volume is $tv/(d+1)$. Therefore
\[
\vol\Pi(\mathcal P K)=(t/d)^d(vH+L/d),\qquad
a(K)=\frac L{dvH}. \tag{2}
\]
Translation replaces $b_i$ by $b_i-u_i\cdot z$, an elementary row operation,
so (2) is translation invariant even when some uncentered $b_i$ vanish.
For the centered cone law, determinant multilinearity gives
$A=d!H/(dv)^d$ and $B=(d+1)!L/(dv)^{d+1}$, proving the identity.
General bodies follow by polytope approximation: choose a common interior
ball and common exterior ball; support functions converge uniformly and
surface-area measures weakly. The cone laws then converge weakly with uniformly
bounded support, so both determinant integrals converge. Projection bodies and
volumes also converge. Positivity of $A,B$ follows from full span and positive
measure neighborhoods of independent support tuples.

\section{Conditioning-free volume-weighted anchors}
The following statement is independent of convex-body realizability.
\begin{lemma}\label{lem:anchors}
For any centered bounded Borel law with full affine support in $\R^d$, and
any norm with support diameter $\Lambda$, there are affinely independent
support anchors $w_0,\ldots,w_d$ and a Borel assignment $I$ such that
\[
\E\|X-w_{I(X)}\|\le\Lambda\frac{(d+1)(d+2)}2\frac DB.
\]
In particular $D\ge0$.
\end{lemma}
\begin{proof}
For a base $\mathbf x=(x_1,\ldots,x_d)$ put
$F_{\mathbf x}(y)=\det((x_1,1),\ldots,(x_d,1),(y,1))$,
$P(\mathbf x)=\E(F_{\mathbf x}(X))_+$ and
$N(\mathbf x)=\E(-F_{\mathbf x}(X))_+$. Affinity and centering give
$D=2\E\min(P,N)\ge0$, including singular bases. Define
\[
\psi_{\mathbf x}(y,z)=\min((F(y))_+,(-F(z))_+)
                       +\min((-F(y))_+,(F(z))_+).
\]
For independent base and test points, each minimum has conditional mean at
most $\min(P,N)$, so $\E\psi\le D$.
For an anchor tuple $W$ write $V(W)=|\det((w_i,1)_{i=0}^d)|$. Form $G(W,x)$
by summing the witness for each base of all anchors except $w_i$, tested at
$w_i,x$, and for each $i<j$ the base of $x$ and all anchors except $w_i,w_j$,
tested at $w_i,w_j$. Orientations are irrelevant. There are
$N_d=(d+1)(d+2)/2$ witnesses, each on $d+2$ independent samples. Thus for
$H(W)=\int G(W,x)d\nu(x)$ we have $\E H\le N_dD$.

When $V>0$ let $\alpha_i(x)$ be anchor barycentric coordinates. The singleton
witness is $V\min(1,(-\alpha_i)_+)$. The pair witness is
$V[\min((\alpha_i)_+,(\alpha_j)_+)+\min((-\alpha_i)_+,(-\alpha_j)_+)]$.
Choose the least index $r$ maximizing $\alpha_i$. Put
$N_x=\sum_i(-\alpha_i)_+$ and $U_x=\sum_{i\ne r}(\alpha_i)_+$.
Affinity and the independent diameter bound give
\[
\begin{aligned}
\|x-w_r\|&\le\Lambda\min(1,N_x+U_x)\\
&\le\Lambda\left[\sum_i\min(1,(-\alpha_i)_+)
                  +\sum_{i<j}\min((\alpha_i)_+,(\alpha_j)_+)\right]\\
&\le\Lambda G(W,x)/V(W).
\end{aligned}
\]
The middle inequality uses the pairs containing the largest positive
coordinate to pay for $U_x$; the capped negative terms pay for
$\min(1,N_x)$. This remains valid for arbitrarily large negative coordinates.

Full affine support gives $B=\E V>0$. Under the probability measure
$d\widehat\mu=(V/B)d\nu^{d+1}$, concentrated on $V>0$,
\[
\E_{\widehat\mu}(H/V)=\E[H\mathbf1_{V>0}]/B\le N_dD/B.
\]
There exists a tuple in this probability-one domain with $H/V$ at most its
mean; otherwise the positive difference has positive integral. Selecting it
and integrating the pointwise bound proves the lemma. This is not a minimum
argument for a possibly singular ratio. The least-index assignment is Borel.
If $D=0$, its cost is zero and the law is supported on its $d+1$ anchors;
their positive probabilities and centering put zero inside their simplex.
\end{proof}

\section{Intrinsic conversion to an enclosing simplex}
We now translate the centroid of the prescribed maximum $S$ to zero, and use
this centering for the cone law. Write $n=d+1$, $Q=K-K$, $N(y)=h_Q(y)$,
whose dual norm on $u$ is the gauge $p_Q(u)$. Replacement maximality gives
$-K\subset dS\subset dK$: the $S$-barycentric coordinates of $-x/d$ are
$(1-\alpha_i(x))/d$. For a cone support point, $h_K(x)=1$, hence
$N(x)\le n$, and the support diameter is at most $2n$.
Lemma~\ref{lem:anchors} gives anchors and assignment of cost
\[
h=\E N(X-w_I)\le\frac{n^3(n+1)\ee}{1+n\ee}\le C\ee,
\qquad C=(d+1)^3(d+2).
\]
For nonzero $u$, every chord of $K$ parallel to $u$ has length at most the
radial function of $Q$ in that direction. Fiber integration and Cauchy's
formula imply, with the homogeneous projection function $\pi_K$,
\[
\pi_K(u)/v\ge p_Q(u),\qquad \E(u\cdot X)_+=\pi_K(u)/(dv)\ge p_Q(u)/d.
\]
Assignment error in that direction is at most $h p_Q(u)$, so
$T=\conv(w_i)$ contains $(1/d-h)Q^\circ$.
When $h<1/(2d)$, $0\in\operatorname{int}T$ and its polar simplex $P$ satisfies
$K\subset P\subset2dQ$. Each $w_i$ is on both polar boundaries and
$h_K(w_i)=h_P(w_i)=1$.

Let $q_i$ be the vertex of $P$ opposite the facet with normal $w_i$, and
$\lambda_i>0$ the barycentric coordinates of zero in $T$. The cone law of $P$
has weights $\lambda_i$: they are its unique centered probabilities. Its
anchor barycentric functions are $\lambda_i(1-q_i\cdot x)$.
Set $p_i=\Pr(I=i)$ and $c=\sum p_iw_i$. Then $N(c)\le h$ and
\[
p_i-\lambda_i=-\lambda_iq_i\cdot c,\qquad
|q_i\cdot c|\le2dh=:r,qquad p_i\ge(1-r)\lambda_i.
\]
For $f_K(u)=\E|u\cdot X|$ and the analogous $f_P$ we obtain
\[
(1-r)f_P(u)\le\E|u\cdot w_I|
\le f_K(u)+h p_Q(u)\le(1+dh/2)f_K(u).
\]
Thus $[\pi_P(u)/\vol P]/[\pi_K(u)/v]\le(1+r/4)/(1-r)$.
Furthermore $z=\E[h_P(X)-1]\in[0,r]$, because
$h_P(y)\le2dN(y)$ and $h_P(w_I)=1$. The mixed-volume formula gives
$1+z=V(K[d-1],P)/v$, and Minkowski's first inequality yields
$\vol P/v\le(1+z)^d\le(1+r)^d$. Consequently
\[
\delta:=\sup_{u\ne0}\frac{\pi_P(u)-\pi_K(u)}{\pi_P(u)}
\le1-\frac{1-r}{(1+r)^d(1+r/4)}\le(d+5/4)r\le3d^2h.
\]
The middle bound follows from $1-abc\le(1-a)+(1-b)+(1-c)$ for factors in
$[0,1]$, and the derivative bounds for $(1+r)^{-d}$ and $(1+r/4)^{-1}$.
Put $H=3d^2C$. If $\ee\le1/H$, then
$h\le1/(3d^2)<1/(2d)$, $r<1$ and $\delta\le H\ee\le1$;
every denominator and strict gate above remains valid at the endpoint.

\section{Relative projection caps retain every maximum simplex}
\begin{lemma}\label{lem:caps}
If $K\subset P$ and $P$ is a simplex, put
$\rho=[\sup_u(\pi_P(u)-\pi_K(u))/\pi_P(u)]^{1/(d-1)}$.
Then every maximum $S\subset K$ satisfies $E(K,S)\le16n^2\rho$.
If $\rho\le1/(16n)$, the stronger bound $E(K,S)\le8n\rho$ holds.
\end{lemma}
\begin{proof}
For each vertex of $P$, a projection direction parallel to its opposite facet
makes the corresponding barycentric coordinate descend to the projection.
If $\max_K\lambda_i=1-t$, the missing projected cap has relative volume
$t^{d-1}$. Hence choose $k_i\in K$ with $\lambda_i(k_i)\ge1-\rho$.
Let $Q_0$ be their nonnegative column-stochastic barycentric matrix.
Independent row choices from its columns show
$\det Q_0\ge2\prod_i(Q_0)_{ii}-1\ge1-2n\rho$.

For $\rho\le1/(16n)$ put $\eta=2n\rho\le1/8$. If $W$ is the barycentric
matrix of the prescribed maximum $S$, maximality gives
$|\det W|\ge1-\eta$. Hadamard's inequality forces every column to have norm
at least $1-\eta$, hence an entry at least $(1-\eta)^2\ge1-2\eta$.
These entries have distinct rows: otherwise two columns collide with
probability at least $(1-2\eta)^2$, and
$|\det W|\le\operatorname{per}W\le4\eta<1-\eta$, a contradiction.
Relabel so $W_{ii}\ge1-4n\rho$. Then
$\|W-I\|_1\le8n\rho\le1/2$. For $B_0=W^{-1}$,
\[
\|B_0-I\|_1\le\frac{8n\rho}{1-8n\rho},\qquad
\beta:=\max_{ij}|(B_0)_{ij}-\mathbf1_{i=j}|
\le\frac{4n\rho}{1-8n\rho}\le1/2.
\]
The factor $1/2$ follows because columns of $B_0-I$ have sum zero.
Put $r_i=1-W_{ii}$. The $ii$ entry of $B_0W=I$ gives
$(B_0)_{ii}-1\ge(1-\beta)r_i/(1-r_i)\ge r_i/2$.
Replacing the $i$th vertex of $S$ by $k_i$ cannot increase volume, so
$|(B_0Q_0)_{ii}|\le1$. Its lower bound is
$(1-\rho)(B_0)_{ii}-\rho\beta$. Therefore
$(B_0)_{ii}-1\le\rho((B_0)_{ii}+\beta)\le2\rho$ and $r_i\le4\rho$.
Now $\|W-I\|_1\le8\rho$, so every negative entry of $B_0$ has magnitude
at most $4\rho/(1-8\rho)\le8\rho$. These columns are precisely the
$S$-barycentric coordinates of the vertices of $P$. A centroid homothety of
excess $t$ is characterized by all barycentric coordinates being at least
$-t/n$. Thus $E(K,S)\le8n\rho$.
For $\rho>1/(16n)$ the universal bound $E\le n$ gives $E<16n^2\rho$.
This proves both assertions without replacing the prescribed maximum $S$.
\end{proof}

\section{Completion of the upper bound and dimension dependence}
On the branch $\ee\le1/H$, the enclosing simplex and
Lemma~\ref{lem:caps} give $E\le16n^2(H\ee)^{1/(d-1)}$.
On the other branch this same right side exceeds $n\ge E$.
This proves the upper bound, including $\ee=0$.
For the explicit coefficient, $H_3=8640<144^2$, and
\[
\frac{H_{d+1}}{H_d}
=\left(\frac{d+1}{d}\right)^2
 \left(\frac{d+2}{d+1}\right)^3\frac{d+3}{d+2}
\le\frac{25}{6}<144\quad(d\ge3).
\]
Induction gives $H_d^{1/(d-1)}\le144$; since $n/d\le4/3$, $G_d\le4096d^2$.
Also $G_d\sim16d^2$. The local linear bound has an exponentially smaller
deficit gate $H\ee\le(16n)^{-(d-1)}$; it does not imply a global $O(d)$
coefficient.

\appendix
\section{The truncation obstruction: complete calculation}

Let $e_1,\ldots,e_d$ be the coordinate vectors and, for $0<t<1$, set
\[
 K_t=\{x\in\mathbb R^d:x_i\ge0,\ t\le x_1+\cdots+x_d\le1\}
     =\conv\{e_1,\ldots,e_d,te_1,\ldots,te_d\}.
 \tag{3}
\]
This is a standard simplex with a homothetic cap cut from one vertex.

\begin{theorem}\label{thm:exact}
For every integer $d\ge2$ and $0<t<1$, the following assertions hold.

\medskip\noindent
\textup{(i)} With $q=t^{d-1}$, the invariant is exactly
\[
 a(K_t)=
 \frac{1+(d-1)q-(d-1)t^d-qt^d}
 {(1-t^d)(d+1+(d-1)q)},
 \tag{4}
\]
and consequently
\[
 \ee(K_t)=
 \frac{t^{d-1}\bigl[d(d-1)-(d+1)(d-2)t-2t^d\bigr]}
 {(d+1)(1-t^d)(d+1+(d-1)t^{d-1})}.
 \tag{5}
\]
In particular,
\[
 \lim_{t\downarrow0}\frac{\ee(K_t)}{t^{d-1}}
   =c_d:=\frac{d(d-1)}{(d+1)^2}>0.
 \tag{6}
\]

\medskip\noindent
\textup{(ii)} The maximum simplex volume is
\[
 M(K_t)=\frac{1-t}{d!}.
 \tag{7}
\]
Every maximum-volume inscribed simplex, including those with vertices in the
relative interiors of faces, has the form
\[
 S_p=\conv\{e_1,\ldots,e_d,tp\},\qquad
 p_i\ge0,\quad\sum_{i=1}^d p_i=1.
 \tag{8}
\]
For each such simplex,
\[
 E(K_t,S_p)=(d+1)t\max_i p_i.
 \tag{9}
\]
Thus
\[
 \inf_{S\text{ maximum}}E(K_t,S)=\frac{d+1}{d}t,
 \qquad
 \sup_{S\text{ maximum}}E(K_t,S)=(d+1)t.
 \tag{10}
\]

\medskip\noindent
\textup{(iii)} The distance to any simplex satisfies
\[
 \left(1+t+\cdots+t^{d-1}\right)^{1/d}
 \le \BM(K_t,\Delta_d)\le 1+t.
 \tag{11}
\]
In fact, for every fixed maximum simplex $S_p$, its smallest enclosing
\emph{translated} homothet has factor exactly $1+t$.
In particular, for $0<t<1$,
\[
 \frac{t}{2d}\le \BM(K_t,\Delta_d)-1\le t,
 \qquad
 \liminf_{t\downarrow0}\frac{\BM(K_t,\Delta_d)-1}{t}\ge\frac1d.
 \tag{12}
\]

\medskip\noindent
\textup{(iv)} The Minkowski asymmetry is exactly
\[
 s(K_t):=\inf\{r\ge1:\exists z,
              -(K_t-z)\subset r(K_t-z)\}=d-t.
 \tag{13}
\]
\end{theorem}

\begin{proof}
We prove each assertion directly.

\medskip\noindent\textbf{Facet determinants.}
Put $c=1/(d-1)!$ and $r=1-q$. The $d$ coordinate facets have area
$c(1-t^{d-1})$. The top and bottom facets have areas $c\sqrt d$ and
$c\sqrt d\,t^{d-1}$, respectively. The facet data in the coordinates (3) are
therefore
\[
 (u_i,b_i)=(-cr e_i,0)\quad(1\le i\le d),\qquad
 (u_+,b_+)=(c\mathbf1,c),\qquad
 (u_-,b_-)=(-cq\mathbf1,-ct^d).
 \tag{14}
\]
The volume is $|K_t|=(1-t^d)/d!$, so $d|K_t|=c(1-t^d)$.

For horizontal $d$-minors, the all-coordinate subset contributes $c^d r^d$;
the $d$ subsets with the top facet contribute $c^d r^{d-1}$ each; and the
$d$ subsets with the bottom facet contribute $c^d q r^{d-1}$ each.
Subsets with both top and bottom have determinant zero because their normals
are parallel. Hence the horizontal sum is
\[
 H=c^d r^{d-1}\bigl[d+1+(d-1)q\bigr].
 \tag{15}
\]
For the lifted $(d+1)$-minors, the subset containing all coordinate facets
and the top facet contributes $c^{d+1}r^d$. Replacing the top facet by the
bottom facet gives $c^{d+1}r^dt^d$. Each of the $d$ subsets with top, bottom and $d-1$
coordinate facets contributes
$c^{d+1}r^{d-1}(q-t^d)$; here $q-t^d=q(1-t)>0$.
Thus the lifted sum is
\[
 L=c^{d+1}r^{d-1}
   \bigl[1+(d-1)q-(d-1)t^d-qt^d\bigr].
 \tag{16}
\]
Dividing $L$ by $d|K_t|H$ gives (4); subtracting $1/(d+1)$ gives (5).
The bracket in (5) decreases strictly to zero as $t\uparrow1$, so the defect
is positive. Dividing (5) by $t^{d-1}$ and passing to $t=0$ proves the exact
limit (6), without any numerical fitting.

\medskip\noindent\textbf{Maximum simplices, including nonvertex ones.}
First consider a simplex selected from the $2d$ vertices in (3). If a ray
$\mathbb R_+e_i$ is unrepresented, the simplex lies in the coordinate
hyperplane $x_i=0$ and has zero volume. Otherwise $d+1$ distinct vertices
on $d$ rays give exactly one doubled ray, with vertices $e_j,te_j$, and one
vertex on every other ray. Subtract the two doubled columns in the lifted
vertex determinant. If $k$ of the other $d-1$ selected vertices are bottom
vertices, its absolute determinant is
\[
 (1-t)t^k.
 \tag{17}
\]
The largest value is $1-t$, attained precisely by all $d$ top vertices and
one bottom vertex. Repeated vertices give determinant zero.

Now let $z_0,\ldots,z_d$ be arbitrary vertices of an inscribed simplex.
Represent each $z_j$ as a convex combination of the $2d$ vertices, and
independently sample a vertex $Y_j$ from that combination. Multiaffinity of
the lifted determinant $D$ gives
\[
 |D(z_0,\ldots,z_d)|
 =|\mathbb E D(Y_0,\ldots,Y_d)|
 \le\mathbb E|D(Y_0,\ldots,Y_d)|\le1-t.
 \tag{18}
\]
This proves the upper bound (7), already attained by the preceding examples.
If the simplex is maximum, equality in (18) implies that
$|D(Y_0,\ldots,Y_d)|=1-t$ almost surely. By (17), exactly one of the $Y_j$
is a bottom vertex and the other $d$ are the $d$ distinct top vertices,
almost surely.

Let $I_j$ be the indicator that $Y_j$ is a bottom vertex. Independence and
$\sum_j I_j=1$ almost surely imply
\[
 0=\operatorname{Var}\Bigl(\sum_j I_j\Bigr)
  =\sum_j \mathbb P(I_j=1)\bigl[1-\mathbb P(I_j=1)\bigr].
\]
Thus every slot has deterministic top/bottom type, with precisely one bottom
slot. The supports of the $d$ top slots must be pairwise disjoint; otherwise
independent sampling gives a repeated top vertex with positive probability.
There are $d$ nonempty disjoint supports in a set of $d$ top vertices, so
each is a singleton. The bottom slot is any convex combination $tp$.
This proves (8). Conversely, the height above the hyperplane
$\sum x_i=1$ of any point $tp$ is fixed, so every simplex (8) has volume
$(1-t)/d!$.

\medskip\noindent\textbf{Centroid containment.}
Write $s_x=\sum_i x_i$. Barycentric coordinates for $S_p$, with index $0$
for $tp$, are
\[
 \beta_0(x)=\frac{1-s_x}{1-t},\qquad
 \beta_i(x)=x_i-tp_i\beta_0(x)\quad(1\le i\le d).
 \tag{19}
\]
Evaluation at the $2d$ vertices of $K_t$ shows
\[
 \min_{K_t}\beta_0=0,\qquad \min_{K_t}\beta_i=-tp_i.
 \tag{20}
\]
For the latter minimum, use $te_j$ with $j\ne i$, available since $d\ge2$.
Affine functions attain their extrema on these vertices.
The centroid-dilated simplex $g_{S_p}+(1+\epsilon)(S_p-g_{S_p})$
is exactly the region $\beta_i\ge-\epsilon/(d+1)$ for every
$0\le i\le d$. Equations (20) give (9), and optimizing $\max_i p_i$
on the probability simplex gives (10).

\medskip\noindent\textbf{Translated containment and distance to any simplex.}
For barycentric coordinates of any simplex $S$, a translated homothet
$v+\lambda S$ has inequalities $\beta_i\ge a_i$ with
$\sum_i a_i=1-\lambda$. Conversely any such thresholds describe a
translated homothet. To contain $K$, one needs
$a_i\le\min_K\beta_i$ for each $i$, so
\[
 \lambda\ge1-\sum_i\min_K\beta_i.
 \tag{21}
\]
Choosing all thresholds equal to their minima realizes equality whenever
$S\subset K$. For (20), the sum is $-t$, giving factor exactly $1+t$.
Concretely,
\[
 K_t\subset -tp+(1+t)S_p.
 \tag{22}
\]
This proves the upper bound in (11) under the full freedom in (1).

For the lower bound, any simplex $S\subset K_t$ has $|S|\le M(K_t)$.
If $K_t\subset v+\lambda S$, volume monotonicity and (7) imply
\[
 \lambda^d\ge\frac{|K_t|}{|S|}
 \ge\frac{|K_t|}{M(K_t)}
 =\frac{1-t^d}{1-t}=1+t+\cdots+t^{d-1}.
 \tag{23}
\]
Taking the infimum proves the lower bound in (11). Furthermore
\[
 (1+t)^{1/d}-1
 =\frac1d\int_0^t(1+u)^{-(d-1)/d}\,du
 \ge\frac{t}{d\,2^{(d-1)/d}}\ge\frac{t}{2d}.
\]
Finally the derivative at zero of
$(1+t+\cdots+t^{d-1})^{1/d}$ is $1/d$, proving (12).
This lower bound concerns \emph{all} simplices, not only maximum simplices.
We do not assert an exact formula for the unrestricted distance in (11).

\medskip\noindent\textbf{Minkowski asymmetry.}
The containment in (13) is equivalent to
$-K_t+(r+1)z\subset rK_t$. Writing $Z=\sum_i z_i$, its necessary and
sufficient facet inequalities are
\[
 (r+1)z_i\ge1\quad(1\le i\le d),\qquad
 1+rt\le(r+1)Z\le r+t.
 \tag{24}
\]
Summing the coordinate inequalities and using the last upper bound gives
$r\ge d-t$. This value is attainable: take
$z_i=1/(d+1-t)$. The bottom inequality follows from
\[
 d-1-(d-t)t=(1-t)(d-1-t)\ge0.
\]
All other inequalities in (24) are equalities. This proves (13).
\end{proof}

\section{Universal implications and necessary constants}

\begin{corollary}\label{cor:power}
Fix $d\ge3$. A power estimate valid for every full-dimensional body of
dimension $d$, even restricted to a sufficiently small lower-end defect,
can have exponent $\alpha$ only if $\alpha\le1/(d-1)$, for each of:
\[
 \BM(K,\Delta_d)-1\le C_d\ee(K)^\alpha;
 \tag{25}
\]
\[
 \inf_{S\text{ maximum}}E(K,S)\le C_d\ee(K)^\alpha;
 \tag{26}
\]
\[
 E(K,S)\le G_d\ee(K)^\alpha
 \quad\text{for every maximum-volume inscribed simplex }S.
 \tag{27}
\]
At the critical exponent $1/(d-1)$, any constant in (27) must satisfy
\[
 G_d\ge(d+1)\left(\frac{(d+1)^2}{d(d-1)}\right)^{1/(d-1)}>d+1.
 \tag{28}
\]
Any constant in (26) must be at least $1/d$ times the right-hand side of
(28), and any constant in (25) must be at least
\[
 \frac1d\left(\frac{(d+1)^2}{d(d-1)}\right)^{1/(d-1)}.
 \tag{29}
\]
These are necessary constants, not claimed optimal universal constants.
In particular a dimension-independent constant is impossible in (27)
at the critical exponent.
\end{corollary}

\begin{proof}
Equations (6), (10), and (12) give a positive constant times $t$ on the
left of each estimate and a constant times $t^{(d-1)\alpha}$ on its right.
If $(d-1)\alpha>1$, division by $t$ and passage to $t\downarrow0$ is a
contradiction. At equality of exponents, take $p=e_1$ in (9) for (28),
$p_i=1/d$ for (26), and the lower-limit bound in (12) for (29).
\end{proof}


\section{Provenance, formal scope and prior work}
The invariant is taken from mxym/math entry 005, whose facet and cone identities
are derived above. The volume-weighted anchor argument, intrinsic conversion
and relative cap bridge consolidate the later research packages; the full
source map and fixed hashes accompany this manuscript. The truncation appendix
retains the original complete computation and all-maximum-simplex classification.
Together they prove the remaining sharpness assertions in
Theorem~\ref{thm:main}. Its formulas are valid also for $d=2$, while the main
upper theorem is stated for $d\ge3$.

The frozen releases
\nolinkurl{formalizations/sharp-simplex-upper-bound} and
\nolinkurl{formalizations/simplex-truncation-sharpness} provide Lean proofs
of the earlier sharp-exponent bound and the truncation obstruction, respectively.
The former has a larger coefficient. They must not be described as Lean
proofs of the present $G_d\le4096d^2$ refinement. Fresh compilation, target
axiom checks and empty-environment kernel replay have distinct verification
records; arithmetic regressions do not prove this analytic refinement.
Full formalization of the improved coefficient and optimal dimension dependence
remain open tasks. No external human review or literature-wide priority
certification is claimed.

Classical tools used above are Cauchy's projection formula, zonotope volumes,
the surface-area and mixed-volume formulas, Minkowski's first inequality,
polytope approximation and Hadamard's inequality. Standard references are
R. Schneider, \emph{Convex Bodies: The Brunn--Minkowski Theory}, 2nd edition,
Cambridge University Press, 2014, and P. McMullen,
\emph{On zonotopes}, Transactions of the AMS 159 (1971), 91--109.
These are inputs, not project results.

Related results include K. J. B\"or\"oczky Jr.,
\emph{The stability of the Rogers--Shephard inequality and of some related
inequalities}, Advances in Mathematics 190 (2005), 47--76,
DOI \nolinkurl{10.1016/j.aim.2003.11.015}. In the plane, the invariant identifies
with a difference-body ratio and the published result gives a direct linear
Banach--Mazur comparison. It does not directly give the prescribed-centroid,
all-maximizer assertion in dimensions $d\ge3$. The accompanying limited
comparison also records cone-volume and random-simplex stability papers.
No absence of an exact match in that screen is treated as proof of novelty.
The exact deficit, prescribed centroid and universal quantifier over maximum
simplices are essential when comparing statements.
\end{document}
